\section{Structural exclusions in dimensions forty and one hundred sixty}
\subsection{A Morita obstruction in dimension forty}
The two rings $R10\_1\_72$ and $R10\_1\_127$ have dimension profile
$(1^4,2^4,\sqrt{10}^{\,2})$. Their invertibles form $V_4$. There is an
involution $h$ fixing each dimension-two simple on both sides, with exact
stabilizer $\{1,h\}$, and exchanging the two irrational-dimensional simples.
Some dimension-two simple is not self-dual. These conditions alone give
the following exclusion, without a pivotal assumption.

Let $\mathcal D$ be the integral subcategory, of dimension twenty.
Since $h$ fixes a simple, its pointed associator is trivial in cohomology:
the fixed simple gives a rank-one module category over $\langle h\rangle$.
Thus $A=1+h$ is a connected separable algebra. There are ten simple right
$A$-modules in $\mathcal D$, all of underlying dimension two: two free
orbits of invertibles, and two module structures on each fixed simple.
For each such module $M$, its internal endomorphism object has dimension
two and contains the unit once. It is therefore $1+g$ for an invertible
$g\ne1$ fixing $M$. The underlying stabilizers force $g=h$.
Every module consequently has two simple left $A$-structures, and
$ {}_A\mathcal D_A$ is pointed of dimension twenty.

Write this pointed category as $\Vec(G,\eta)$. The Morita module has
rank ten, so $\mathcal D$ is a group-theoretical category
$\mathcal C(G,\eta,K,\psi)$ with $|K|=2$. The groups of order twenty
are the two abelian groups, $D_{20}$, $\operatorname{Dic}_5$, and
$F_{20}=C_5\rtimes C_4$ with faithful action. A normal order-two $K$
would make every simple invertible, excluding the abelian and dicyclic
possibilities. In $D_{20}$ a nonnormal $K$ is a reflection subgroup;
every relevant double coset is inverse-stable with trivial stabilizer,
so every dimension-two simple is self-dual, independently of $\eta,\psi$.
Hence $G=F_{20}$.

The full Morita dual $ {}_A\mathcal C_A$ has twenty invertibles and two
simples of dimension $\sqrt{10}$. Its noninvertible stabilizer has order
ten and must be abelian by Liptrap's generalized Tambara--Yamagami theorem.
But $F_{20}$ has a unique index-two subgroup, the nonabelian $D_{10}$.
This contradiction proves both exclusions. The proof allows arbitrary
associators throughout; it does not assume that a known untwisted model
is the only possible categorification of the integral subring.

\subsection{Forced commutativity in dimension one hundred sixty}
The two rings $R10\_1\_83$ and $R10\_1\_177$ have dimensions
$(1^2,2^2,5^6)$ and trivial universal grading. Their nontrivial invertible
$g$ fixes the two dimension-two simples and exchanges the six dimension-five
simples in pairs. ENO's categorical Burnside theorem makes a hypothetical
category solvable. In a prime cyclic extension/equivariantization chain,
the final operation cannot be an extension, because the universal grading
is trivial. It must be a $C_2$-equivariantization
$\mathcal C=\mathcal D^{C_2}$.

The sign permutation forces $\mathcal D$ to have dimensions $(1^5,5^3)$.
Its invertibles form $C_5$, act trivially on the three noninvertibles, and
belong to $\mathcal D_{\rm ad}$. The possible adjoint dimensions are
$5+25k$ and must divide eighty; only five and eighty qualify. Dimension
five cannot contain a dimension-five simple in any homogeneous component,
so $\mathcal D_{\rm ad}=\mathcal D$. A second application of the solvable
last-step argument gives a central $\Rep(C_5)$ whose image is exactly the
five invertibles of $\mathcal D$.

An invertible object has at most one half-braiding when the universal grading
is trivial: the ratio of two half-braidings is a tensor automorphism of
the identity, hence a character of that grading. This central $\Rep(C_5)$
is therefore stable under the $C_2$ action. Its equivariantization is a
central Tannakian subcategory of $\mathcal C$ with dimensions $(1,1,2,2)$,
namely $\Rep(D_{10})$. Direct de-equivariantization gives
$\mathcal C=\mathcal E^{D_{10}}$ with $\FPdim\mathcal E=16$.
The orbit-stabilizer dimension formula forces $\mathcal E$ to be pointed:
every nonunit orbit contributes a dimension-five simple and hence has five
invertible underlying objects. The order-five rotation acts freely on all
fifteen nonidentity elements of its group. Characteristic-center and
two-torsion arguments force that group to be $V=C_2^4$.

Identify $V$ with $\mathbb F_{16}$. A rotation is multiplication by an
order-five element $a$, and a reflection can be taken to be $b(x)=x^4$.
The complete presentation of $H^3(V,\C^*)$ has fourteen binary generators.
Their restrictions to cyclic subgroups are the squarefree monomials of
degrees one through three on $V\setminus\{0\}$. They span the even-parity
subspace of its fifteen functions. Rotation invariance leaves four classes,
with constant values on the three length-five orbits and even total parity.
Four explicit normal forms in the certificate cover these classes.

The rotation tensorator can be made strict because
$H^1(C_5,H^2(V,\C^*))=0$. Relative to a reference reflection tensorator,
the commutator of any other reflection tensorator is a rotation-invariant
alternating form. Every such form is
\[
 B_c(x,y)=\operatorname{Tr}_{16/2}(cxy^4),\qquad c\in\mathbb F_4.
\]
It vanishes on $\mathbb F_4\times\mathbb F_4$, the reflection's fixed
space. The reference tensorators are also symmetric there, as checked
exactly. Consequently every allowed tensorator is symmetric on fixed pairs.
This argument covers arbitrary complex actions, not only a bounded search
of roots of unity.

For equivariant objects $X,Y$, the trace of an element $h$ on their tensor
product's fixed homogeneous degree $z$ is
\[
 \sum_{\substack{x+y=z\\hx=x,\ hy=y}}
 J_h(x,y)^{-1}\chi_X(h,x)\chi_Y(h,y).
\]
The established symmetry makes this invariant under exchanging $X,Y$.
Projective stabilizer characters determine equivariant objects, so the
Grothendieck ring must be commutative. Both input tensors are explicitly
noncommutative, proving the exclusions. The full proof and independent
standard-library replay are in \nolinkurl{work/audit_new_models/INT160_AUDIT.md}
and \nolinkurl{verify_int160_audit.py}.

\section{The generalized Tambara--Yamagami family}
Fourteen rank-ten inputs have profile $(1^8,2^2)$. Liptrap's classification
of the corresponding fusion systems supplies necessary and sufficient
associator data. Eleven have explicit unitary data in the package. The
remaining three, $R10\_1\_118$, $R10\_1\_122$, and $R10\_1\_180$, admit no
complex solution of those equations.

For a positive record the phases are powers of a sixteenth root of unity,
with $\tau=1/2$. The verifier checks the defining biderivation and coherence
equations and all eighty-eight recorded nontrivial Fourier blocks for
unitarity. Liptrap's construction then supplies the entire fusion category.
For a negative record, the stabilizer bicharacters are enumerated completely
(sixteen pairs for $V_4$, four for $C_4$). The remaining multiplicative
equations have the form $\prod_jx_j^{A_{ij}}=i^{b_i}$ with arbitrary
$x_j\in\C^*$. For every allowed bicharacter choice the receipt gives an
integer row vector $y$ with
\[
 yA=0,\qquad yb\not\equiv0\pmod4.
\]
Multiplying the equations with these integer exponents would imply
$1=i^{yb}\ne1$. The obstruction therefore applies over all of $\C^*$;
failure of a finite-phase search is not being used as nonexistence evidence.
